\subsection{慢性腰痛与腰椎问题者的性姿势}

腰椎间盘突出、腰肌劳损、坐骨神经痛是成年人最常见的几类腰背问题，它们\textbf{并不禁止性生活}——处于缓解期的腰痛者完全可以维持满意的亲密，需要的只是姿势上的调整。需要暂缓的情况是：急性发作期（剧痛、下肢放射痛、麻木无力），先把治疗和休息放在前面。

理解姿势的原则，先理解腰椎最怕什么。\textbf{前屈 + 旋转 + 负重}是腰椎最不喜欢的组合——弯着腰用力、再带一点扭转的动作，正是搬东西闪到腰的典型姿势。性生活里也一样：凡是让腰椎离开"中立位"（自然的生理曲度）、又需要腰腹持续发力的姿势，对腰椎问题者都是高风险的。反过来，让脊柱保持自然曲线、由大肌群而非腰部代偿发力的姿势，就是相对安全的选择。

\begin{itemize}
	\item \textbf{侧卧位（汤匙式）是腰痛者的首选}：双方侧躺，脊柱处于中立位，腰椎负荷最小，动作幅度也容易控制（详见姿势章汤匙式姿势条目）；腰部不适的一侧在上或在下可以自行试验，取疼痛最小的方向。
	\item \textbf{仰卧位 + 膝下垫枕}：平躺时在膝盖下垫一个枕头，让髋膝微屈，可以明显减少腰椎前凸的压力——这个准备动作简单到几乎不应该被省略。
	\item \textbf{女上位可以采用，但要注意}：由骨盆问题较轻或腰部健康的一方主导节奏与深度，通常更安全；采用时上方者用膝支撑而非腰发力，避免腰部过度后仰。
\end{itemize}

相对需要\textbf{慎重}的姿势：需要腰部大幅后伸的深弓背后入式、需要一方手臂持续支撑躯干的传统男上位（腰椎与肩肘同时在代偿）、以及任何做起来"腰部发紧"的旋转动作。一个实用的自检标准是：\textbf{如果某种姿势在日常生活中都会让你腰痛，它没有理由在床上例外}。

几个可以减轻负担的细节：性生活前先热敷腰部 15--20 分钟（与运动前热身同理）；枕头不只垫头——垫在膝下、腰下、两膝之间都各有用处；不必为了性生活临时戴上护腰——护腰会让深层肌肉在关键时刻"偷懒"，日常问题请遵医嘱使用。

如果你做过的不是腰椎手术而是\textbf{髋或膝关节置换}：术后早期人工髋关节的屈曲角度通常不能超过 90 度、且要避开手术侧的受压与内收动作，何时恢复性生活、能采用哪些姿势，请以手术医生的个体化医嘱为准，不要照搬他人经验。

最后是一条不能讨价还价的红线：如果性生活之后出现\textbf{下肢放射痛或麻木加重、会阴区麻木、大小便控制异常}——这些是马尾综合征的警示信号，请立即就医，不要等待观察。性生活可以调整，神经的损伤等不起。腰部问题较重、姿势受限较大的伴侣，还可以参考本章后文"残疾伴侣辅助姿势"一节，其中"减少一方体力负担"的思路完全通用。
